\section*{Capítulo \ref{ch:b2:poynting-vector} --- Energia eletromagnética e o vetor de Poynting}

\begin{solution}{exo:b2:poynting-vector:1}
$E_0 = \sqrt{2I/\varepsilon_0c} = \qty{870}{V/m}$, $B_0 = E_0/c = \qty{2.9}{\micro T}$; $u = I/c
= \qty{3.3}{\micro J/m^3}$; $P = I/c = \qty{3.3}{\micro Pa}$ (telhado preto), \qty{6.7}{\micro Pa}
(espelho); \qty{0.33}{mN} no telhado; $I/hf = 1000/3.2 \times 10^{-19} = \qty{3e21}{m^{-2}.s^{-1}}$.
\end{solution}

\begin{solution}{exo:b2:poynting-vector:2}
Área \qty{1.8}{mm^2}: $I = \qty{2.8}{kW/m^2}$, $E_0 = \qty{1.5}{kV/m}$, $B_0 = \qty{4.9}{\micro T}$.
Focalizado em \qty{5}{\micro m}: $I = \qty{2.5e8}{W/m^2}$, $E_0 = \qty{4.4e5}{V/m}$. Petawatt:
$I = \qty{1e26}{W/m^2}$, $E_0 = \qty{2.7e14}{V/m}$, quinhentas vezes o campo
atômico --- o átomo é despedaçado em um ciclo.
\end{solution}

\begin{solution}{exo:b2:poynting-vector:3}
$\tfrac12\varepsilon_0E^2$: \qty{40}{J/m^3}, \qty{4.4e4}{J/m^3}. $B^2/2\mu_0$: \qty{4}{kJ/m^3},
\qty{0.9}{MJ/m^3}, \qty{8e8}{J/m^3}. O túnel da máquina de ressonância guarda \qty{0.9}{MJ} --- \qty{26}{mL}
de gasolina.
\end{solution}

\begin{solution}{exo:b2:poynting-vector:4}
$F = 2IA/c = \qty{0.09}{N}$; $a = \qty{4.5e-4}{m/s^2}$; \qty{1.2}{km/s} em um mês;
o Sol puxa com \qty{1.2}{N}: a vela não pode vencer a gravidade, mas bordeja
contra ela em órbita.
\end{solution}

\begin{solution}{exo:b2:poynting-vector:5}
(a) $j = \qty{3.2e6}{A/m^2}$, $E = j/\gamma = \qty{0.053}{V/m}$; $B = \mu_0I/2\pi a = \qty{2}{mT}$;
$\Pi = EB/\mu_0 = \qty{84}{W/m^2}$, radialmente para dentro. (b) $84 \times 2\pi a = \qty{0.53}{W}$
por metro; $R = 1/\gamma\pi a^2 = \qty{5.3}{m\ohm/m}$, $RI^2 = \qty{0.53}{W}$. (c) $\Pi(r)
= EIr/2\pi a^2$, de modo que $2\pi r\Pi = EIr^2/a^2$ e a casca recebe $\dd(2\pi r\Pi)
= 2EIr\,\dd r/a^2 = (j^2/\gamma)2\pi r\,\dd r$. (d) As cargas superficiais da pilha
estabelecem um campo elétrico ao longo de todo o circuito, fora dos fios;
o $\vect E\wedge\vect B$ ali leva a energia da pilha até a carga.
\end{solution}

\begin{solution}{exo:b2:poynting-vector:6}
(a) $E = Q/\varepsilon_0\pi R^2$ axial, $B = \mu_0Ir/2\pi R^2$ azimutal, $\vect\Pi = -\varepsilon_0
(r/2)E\dot E\,\vect e_r$. (b) $\varepsilon_0(R/2)E\dot E\cdot2\pi Rd = \dd(\tfrac12\varepsilon_0E^2\pi R^2d)/
\dd t$. (c) $\dot E = I/\varepsilon_0\pi R^2 = \qty{1.4e13}{V.m^{-1}.s^{-1}}$; $\Pi(R) = \qty{3.2e6}{W/m^2}$;
vezes $2\pi Rd = \qty{3.1e-4}{m^2}$: \qty{1.0}{kW} $= UI$ com $U = Ed = \qty{1}{kV}$. (d)
Sim: menor por fatores do tipo $(\omega R/c)^2$ (\cref{ch:b2:maxwell-equations}).
\end{solution}

\begin{solution}{exo:b2:poynting-vector:7}
(a) $\vect E = -(r/2)\dot B\,\vect e_\theta$, $\vect\Pi = -(rB\dot B/2\mu_0)\vect e_r$: para dentro
enquanto $B$ cresce. (b) $(RB\dot B/2\mu_0)2\pi R = \dd(\pi R^2B^2/2\mu_0)/\dd t$. (c)
Fora, $\vect B = \vect 0$, logo $\vect\Pi = \vect 0$: a energia é injetada no
próprio enrolamento, onde o gerador força a corrente contra o
campo induzido (ali $\vect j\cdot\vect E < 0$); logo por dentro da lâmina
$\vect\Pi$ já é o fluxo para dentro de (a).
\end{solution}

\begin{solution}{exo:b2:poynting-vector:8}
(a) $E = U/r\ln(b/a)$ radial, $B = \mu_0I/2\pi r$ azimutal. (b) $\vect\Pi = (UI/
2\pi r^2\ln(b/a))\vect e_z$, no sentido da corrente do núcleo. (c) $\int_a^b\Pi\,2\pi r\,\dd r =
UI$. (d) $E$ ganha uma componente axial $j/\gamma$; $\vect\Pi$ inclina-se rumo ao
núcleo, e seu fluxo radial para dentro de um metro de núcleo vale $RI^2$ por metro.
\end{solution}

\begin{solution}{exo:b2:poynting-vector:9}
(a) $F_{\text{rad}} = I_0(D_0/D)^2\pi a^2/c$, $F_g = GM\cdot\tfrac43\pi a^3\rho/D^2$: razão
$3I_0D_0^2/4cGM\rho a$, independente de $D$. (b) Razão $> 1$ para $a < 3I_0D_0^2/
4cGM\rho = \qty{0.3}{\micro m}$. (c) A poeira fina é empurrada para longe do Sol, mas
guarda a celeridade orbital do cometa, de modo que a cauda de poeira aponta para longe e
se encurva; a cauda de íons, impelida pelo vento solar, é reta.
\end{solution}

\begin{solution}{exo:b2:poynting-vector:10}
(a) $\dot E < 0$: $\vect\Pi$ aponta para fora na borda. (b) $-\dd W_e/\dd t$
sai do capacitor e $RI^2$ entra no fio; no regime quase estático
os dois são iguais --- nada se acumula no meio. (c) $W_e =
Q^2/2C$, $\dd W_e/\dd t = (Q/C)\dd Q/\dd t = -UI = -RI^2$. (d) $\tfrac12LI^2$
anula-se nas duas pontas ($I = 0$) e é desprezível no meio, para uma espira
pequena.
\end{solution}

\begin{solution}{exo:b2:poynting-vector:11}
(a) $\vect\Pi = (EB/\mu_0)\vect e_y$, uniforme: $\operatorname{div}\vect\Pi = 0$. (b) Nulo.
(c) Segundo $y$ pelo vão, fechando-se pelos campos de borda nas
extremidades das placas e do ímã. (d) Um corpo em repouso em campos
estáticos não conduz corrente: $\vect j\cdot\vect E = 0$, e não há aquecimento. O
$\vect\Pi$ circulante é contabilidade, não um fluxo que deposita energia.
\end{solution}

\begin{solution}{exo:b2:poynting-vector:12}
Faraday: $\partial_xE_y = -2E_0k\sin kx\cos\omega t$ é igual a $-\partial_tB_z$, pois
$\omega = kc$. (a)
$\vect\Pi = (E_yB_z/\mu_0)\vect e_x = -(E_0^2/\mu_0c)\sin2kx\sin2\omega t\,\vect e_x$: média
nula. (b) $u = 2\varepsilon_0E_0^2(\cos^2kx\cos^2\omega t + \sin^2kx\sin^2\omega t)$: elétrica
em $x = 0, \lambda/2, \dots$, magnética em $\lambda/4, \dots$, indo e vindo entre elas
duas vezes por período. (c) Quantidade de movimento: cada onda tem $I = \varepsilon_0cE_0^2/2$, e a
refletida inverte a sua: $2I/c = \varepsilon_0E_0^2$. Laplace: em $x =
0$, $B = (2E_0/c)\sin\omega t$ logo do lado de fora e $0$ dentro, de modo que $j_s = B/\mu_0$; a
força por unidade de área sobre a lâmina vale $\tfrac12j_sB = B^2/2\mu_0 = 2\varepsilon_0E_0^2\sin^2
\omega t$. (d) Sua média vale $\varepsilon_0E_0^2 = 2I/c$.
\end{solution}

\begin{solution}{pb:b2:poynting-vector:1}
\textbf{1.} $E = U/r\ln(b/a)$: \qty{1.8}{MV/m} em $r = a$; $B = \mu_0I/2\pi r$:
\qty{10}{mT} em $a$.

\textbf{2.} $\vect\Pi = (UI/2\pi r^2\ln(b/a))\vect e_z$; $\int_a^b\Pi\,2\pi r\,\dd r = UI =
\qty{10}{MW}$.

\textbf{3.} $\tfrac12\varepsilon_0\varepsilon_rE^2 = \qty{34}{J/m^3}$, $B^2/2\mu_0 = \qty{40}{J/m^3}$:
comparáveis.

\textbf{4.} $\Pi(a) = EB/\mu_0 = \qty{1.45e10}{W/m^2}$, $u = \qty{74}{J/m^3}$: $\Pi/u =
\qty{2.0e8}{m/s} = c/\sqrt{\varepsilon_r}$.

\textbf{5.} $E_z = I/\gamma\pi a^2 = \qty{0.027}{V/m}$; $\Pi_r = E_zB/\mu_0 = \qty{210}{W/m^2}$
para dentro; $\times2\pi a = \qty{13}{W/m}$; $R = 1/\gamma\pi a^2 = \qty{53}{\micro\ohm/m}$, $RI^2 =
\qty{13}{W/m}$.

\textbf{6.} $\Pi_r/\Pi_z = E_z/E_r = 1.5 \times 10^{-8}$: quase toda a energia corre
ao longo da linha; um sopro dela vaza para o cobre.

\textbf{7.} \qty{133}{kW} perdidos de \qty{10}{MW}, $1.3\%$; a \qty{200}{kV}: \qty{1.3}{kW},
$0.013\%$.

\textbf{8.} $\int_a^b(B^2/2\mu_0)2\pi r\,\dd r = (\mu_0I^2/4\pi)\ln(b/a) = \qty{27}{mJ/m} =
\tfrac12\Lambda I^2$ com $\Lambda = \mu_0\ln(b/a)/2\pi = \qty{0.22}{\micro H/m}$; elétrica:
$\tfrac12\Gamma U^2$ com $\Gamma = 2\pi\varepsilon_0\varepsilon_r/\ln(b/a) = \qty{117}{pF/m}$: \qty{23}{mJ/m}.

\textbf{9.} $E_0 = \qty{870}{V/m}$, $B_0 = \qty{2.9}{\micro T}$, $u = \qty{3.3}{\micro J/m^3}$,
\qty{3e21}{m^{-2}.s^{-1}}.

\textbf{10.} \qty{320}{W} elétricos; \qty{1280}{W} de calor em \qty{3.2}{m^2} de
faces: $\Delta T = 1280/32 = \qty{40}{K}$.

\textbf{11.} $P = I/c = \qty{3.3}{\micro Pa}$; \qty{5.3}{\micro N} contra um peso de
\qty{180}{N}.

\textbf{12.} $320 \times 6 = \qty{1.9}{kWh/dia}$, \qty{700}{kWh/ano}; cinco painéis,
\qty{8}{m^2}.

\textbf{13.} $30 \times 10.7 = \qty{320}{W}$; a potência flui no campo entre
os dois fios --- $\vect E$ ao longo dos \qty{30}{V}, $\vect B$ em torno das correntes ---
e não no cobre.

\textbf{14.} $4\pi D^2I = \qty{3.9e26}{W}$; $P/c^2 = \qty{4.3e9}{kg/s}$.

\textbf{15.} $5 \times 10^{21} \times 1.6 = \qty{5e21}{}$ fótons por segundo.

\textbf{16.} $I\pi R_E^2/c = \qty{5.8e8}{N}$: $10^{-14}$ da gravidade.

\textbf{17.} $\lambda = \qty{12.2}{cm}$; a cavidade de \qty{30}{cm} tem alguns comprimentos de onda
--- um ressoador com modos, e não um feixe livre; as paredes metálicas refletem a
onda (\omterm{def:b2:dispersion-wave-packets:complex}{profundidade de penetração} de micrômetros) e a retêm.

\textbf{18.} $0.25 \times 4180 \times 60 = \qty{63}{kJ}$ a \qty{640}{W}: \qty{98}{s};
$640/2.5 \times 10^{-4} = \qty{2.6}{MW/m^3}$.

\textbf{19.} $I = 640/5 \times 10^{-3} = \qty{1.3e5}{W/m^2}$, $E_0 = \qty{10}{kV/m}$; pontos quentes
a $\lambda/2 = \qty{6}{cm}$ um do outro --- daí o prato giratório.

\textbf{20.} $W = QP/\omega = 10^3 \times 800/1.54 \times 10^{10} = \qty{52}{\micro J}$; $u =
\qty{1.7}{mJ/m^3}$; $E_0 = \sqrt{2u/\varepsilon_0} = \qty{20}{kV/m}$ --- cem vezes abaixo
da ruptura.

\textbf{21.} $\tfrac12\omega\varepsilon_0\varepsilon''E_0^2 = \qty{2.6e6}{W/m^3}$: $E_0 = \qty{2}{kV/m}$
dentro da água --- da ordem do campo da cavidade, reduzido pela
permissividade da água: coerente.

\textbf{22.} $hf = \qty{1.6e-24}{J} = \qty{1e-5}{eV}$; $4 \times 10^{26}$ fótons por segundo;
um milhão de vezes pouca energia por fóton para ionizar.

\textbf{23.} $\sim u = \qty{2}{mPa}$: não.

\textbf{24.} Sem nada que a absorva, o fluxo de $\vect\Pi$ não tem
onde terminar senão de volta no magnetron, que superaquece.

\textbf{25.} Linha: o dielétrico, $\vect\Pi$ axial, \qty{10}{MW}, uma carga na
ponta. Painel: o feixe, $\vect\Pi$ ao longo dos raios do Sol, \qty{1.6}{kW}
de entrada, \qty{320}{W} de saída, o semicondutor. Forno: o campo da cavidade, $\vect\Pi$
circulante, \qty{800}{W}, os dipolos da água.
\end{solution}
